\begin{trapbox}
	\begin{description}[style=nextline, leftmargin=2.8em, labelsep=0.5em, itemsep=0.35em]
		\item[\textbf{\textcolor{CTrap}{易错点一（条件忽略）}}]应用公式时，未先确认弦$AB$是否与$y$轴平行，或中点纵坐标$y_0$是否为0。
		\item[\textbf{\textcolor{CTrap}{正确理解（条件验证）：}}]必须在$AB$不与$y$轴平行（斜率存在）且$y_0\neq0$的条件下才能使用$k=p/y_0$。
		\item[\textbf{\textcolor{CTrap}{错因分析（死记形式）：}}]只记忆了公式的外形，忽略了公式的推导前提。
	\end{description}

	\medskip
	\begin{description}[style=nextline, leftmargin=2.8em, labelsep=0.5em, itemsep=0.35em]
		\item[\textbf{\textcolor{CTrap}{易错点二（$p$值误判）}}]从抛物线方程$y^2=4x$直接读取$p=4$，导致斜率计算错误。
		\item[\textbf{\textcolor{CTrap}{正确理解（标准对照）：}}]标准方程$y^2=2px$中$x$一次项系数为$2p$，故$p$应取系数的一半。例如$y^2=4x$中$p=2$。
		\item[\textbf{\textcolor{CTrap}{错因分析（生搬硬套）：}}]未将方程化为标准形式对比，盲目代入系数。
	\end{description}

	\medskip
	\begin{description}[style=nextline, leftmargin=2.8em, labelsep=0.5em, itemsep=0.35em]
		\item[\textbf{\textcolor{CTrap}{易错点三（符号忽视）}}]只考虑$p>0$的情形，当$p<0$时仍使用$p$的正值，或者认为$k$总为正。
		\item[\textbf{\textcolor{CTrap}{正确理解（符号一致）：}}]公式中$p$可正可负，$k$的符号由$p$与$y_0$共同决定。
		\item[\textbf{\textcolor{CTrap}{错因分析（思维定势）：}}]习惯性认为$p$为正数，忽略题设中$p\neq0$包含负数可能。
	\end{description}
\end{trapbox}
